\documentclass[10pt,a4paper]{amsart}
\usepackage[margin=19mm]{geometry}
\usepackage{amsmath,amssymb,mathtools}
\usepackage[scr=rsfs,cal=euler]{mathalfa}
\usepackage{xfrac}
\usepackage[unicode,hidelinks,bookmarks=false]{hyperref}
\newcommand{\ie}{i.e.\ }
\newcommand{\cf}{cf.\ }
\newcommand{\resp}{respectively\ }
\newcommand{\Subsection}{\subsection}
\providecommand{\nobreakdash}{\nobreak\hbox{-}}
\setlength{\emergencystretch}{2em}
\multlinegap0pt
\usepackage{mathtools}
\usepackage{amssymb}
\usepackage{algorithm}\usepackage{algpseudocode}
\usepackage{xcolor}
\newtheorem{lemma}{Lemma}
\newcommand{\GW}{\mathrm{GW}}
\newcommand{\Leb}{\mathrm{Leb}}
\renewcommand{\P}{\mathcal{P}}
\newcommand{\R}{\mathbb{R}}
\newcommand{\X}{\mathcal{X}}
\newcommand{\Y}{\mathcal{Y}}
\title{Beyond Procrustes distances: a multilinear Gromov-Wasserstein distance capturing chirality}
\author{Cl\'ement Soubrier, Geoffrey Woollard, Andrew Warren, Khanh Dao Duc}
\date{Source: arXiv:2608.27774v1 (CC BY 4.0)}
\begin{document}
\maketitle

\noindent\emph{Source and licence.} Beyond Procrustes distances: a multilinear Gromov-Wasserstein distance capturing chirality. Version arXiv:2608.27774v1 (CC BY 4.0), distributed by arXiv under the Creative Commons Attribution 4.0 International licence, http://creativecommons.org/licenses/by/4.0/, as recorded in the arXiv licence metadata for this exact version and shown on the abstract page https://arxiv.org/abs/2608.27774v1. Full licence notice: \texttt{licenses/arxiv-2608.27774-CC-BY-4.0.txt}.\par\smallskip

\noindent\textit{Editorial scope.} Counted unit: the article's lemma lem:Monge\_DGW (source0.tex lines 334--336). The article's complete original proof of it is delivered with it (source0.tex lines 735--755, one \texttt{proof} environment of 21 lines, which the article places away from the statement). No mathematics is written, repaired or re-derived: the statement and the proof are the article's own lines, in the article's own notation, and no retained line is substituted or re-flowed.  Retained uncounted as the source-local context the proof needs: lines 706--711.  Nothing was omitted from any retained span, so the package carries no omission record.  No retained line contains a citation, so the wrapper carries no bibliography and no bibliography entry.  No retained line contains a figure environment, an image-inclusion command or drawing code, so no image is delivered and the package declares no asset dependency.  Editorial formatting only, in addition: an AMS \texttt{amsart} preamble that restates the article's own declarations and macros used by the retained lines, namely the environments \texttt{lemma} on separate counters, together with the standard packages those lines need.  The retained environments print in this document as Lemma 1 in order of appearance, whereas the article prints the same items with the numbering of its own full document.  The complete native source of the article as distributed in the arXiv e-print for this exact version is bundled unchanged as \texttt{sources/208740/source0.tex}.
\begin{lemma}\label{lem:non_dege_small_set}
    Let $A\subset\R^d$ and $m$ be a non-degenerate $n$-form on $\R^d$. Assume that there exists $x_1,\dots,x_d\in A$ a basis of $\R^d$. Then, 
    \[\{T: A\to\R^d, m\circ T^{\times n}=m\}\subset GL(d),\] 
    and the unique linear extension of $T$ to $\R^d$, $\tilde{T}$, verifies $m\circ \tilde{T}^{\times n}=m$.
    
\end{lemma}

\begin{lemma}[Existence of a Monge map for $\GW_m$] \label{lem:Monge_DGW}
    Suppose that $\mu,\nu \in \P(\R^d)$ have compact support and that $\mu\ll \Leb$, i.e. $\mu$ has a density. Then there exists an optimal correspondence map $T:\R^d\to\R^d$. Thus, the optimal transport plan for the $\GW_{m}$ problem can be written as $(id,T)_\sharp\mu$. Moreover, any optimal transport plan is a Monge map.
\end{lemma}

\begin{proof}
    Let $\X\subset\R^d$ be compact and $\mu  \in \P(\X)$ have a density with respect to the Lebesgue measure. Note that this implies that $\X$ has non-zero Lebesgue measure. We just need to prove that for $\nu\in \P(\X)$: \[  \GW_{m}(\mu,\nu)= 0 \implies \exists g\in G, \mu = g_\sharp\nu.\]
    We suppose that $\GW_{m}(\mu,\nu)= 0$. Under the compactness and density assumptions, Lemma \ref{lem:Monge_DGW} proves the existence of an  optimal Monge map $T:\X\to\Y$. In other words, an optimal solution $\pi^*$ of the $\GW_{m}(\mu,\nu)$ problem can be written $\pi^*=(id,T)_\sharp\mu$; and $\nu =T_{\sharp} \mu$. Accordingly,
    \begin{equation*}
        \begin{split}
             0&=\GW_{m}(\mu,\nu)^2\\
             &=\inf_{\pi\in\Pi(\mu,\nu)}\int_{(\X\times\X)^n}(m(\mathbf{x})-m(\mathbf{y}))^2d\pi^{\otimes n}(\mathbf{x},\mathbf{y}) \\
             &=\int_{(\X\times\X)^n}(m(\mathbf{x})-m(\mathbf{y}))^2d(\pi^*)^{\otimes n}(\mathbf{x},\mathbf{y}) \\
             &=\int_{\X^n}(m(\mathbf{x})-m(T^{\times n}(\mathbf{x})))^2d\mu^{n}(\mathbf{x}). \\
        \end{split}
    \end{equation*}
    It remains to show that $T\in GL(d)$. 
    
    Without loss of generality, we can suppose that $m$ is non-degenerate at slot 1.
    Using the multilinearity of $m$, we can write:
    \[\forall x, x_2\dots,x_n\in \R^d,\quad m(x, x_2,\dots,x_n) = \langle x, f(x_2,\dots,x_n)\rangle,\]
    with $f:(\R^d)^{n-1}\to\R^d$ whose components are  multilinear $(n-1)$-forms. We define \[B_1\coloneqq \{\mathbf{x}_1,\dots,\mathbf{x}_{d} \in (\X)^{n-1}, (f(\mathbf{x}_1),\dots,f(\mathbf{x}_d))\text{ is a basis of }\R^d\} ,\] and we have  \[B_1^c  =  \{\mathbf{x}_1,\dots,\mathbf{x}_{d} \in (\X)^{n-1}, \;\det(f(\mathbf{x}_1),\dots,f(\mathbf{x}_d))=0\}.\] But, $\det(f(\mathbf{x}_1),\dots,f(\mathbf{x}_d))=0$ is a polynomial equation, since $f$ has polynomial coordinates. Now the non-degeneracy assumption implies that $B_1$ is non-empty, thus $B_1^c\subsetneq \R^{d^2(n-1)}$. As the set of solutions of a polynomial equation $\Leb(B_1^c)=0$ and $\mu^{d(n-1)}(B_1)=1$. Similarly, the following set $B_2\coloneqq \{x_1,\dots,x_{d} \in (\X)^{d}, (x_1,\dots,x_{d})\text{ is a basis of }\R^d\} $ has measure one. 
    
    Let's fix $A\subset\X$ with $\mu(A)>0$. By modifying $A$ on a set of $\mu$-measure zero, we can assume that $m\circ T^{\times n}=m$ on all of $A$. Observe that $\mu^{dn}((B_2\times B_1)\cap  A^{dn})=\mu(A)>0$.
    By construction, $A$ and $m$ satisfy the assumptions of Lemma \ref{lem:non_dege_small_set}. Hence, there exists $T\in GL(d)\cap G$, such that $\mu = T_\sharp \nu$. In particular, $T\in G$.
\end{proof}

\end{document}
